\documentclass[11pt,a4paper]{article}
\usepackage[utf8]{inputenc}
\usepackage{amsmath,amssymb,amsthm}
\usepackage{listings}
\usepackage{xcolor}
\usepackage[margin=1in]{geometry}
\usepackage{hyperref}

\newtheorem{theorem}{Theorem}

\lstset{
  language=C++,
  basicstyle=\ttfamily\small,
  keywordstyle=\color{blue}\bfseries,
  commentstyle=\color{green!60!black},
  stringstyle=\color{red},
  numbers=left,
  numberstyle=\tiny\color{gray},
  breaklines=true,
  frame=single,
  tabsize=4,
  showstringspaces=false
}

\title{IOI 2001 -- Mobile Phones}
\author{Solution}
\date{}

\begin{document}
\maketitle

\section{Problem Statement Summary}
Maintain an $S \times S$ grid (initially all zeros, $S \le 1024$)
under two operations:
\begin{enumerate}
  \item $\textsc{Update}(x, y, v)$: add $v$ to cell $(x, y)$.
  \item $\textsc{Query}(x_1, y_1, x_2, y_2)$: return the sum of all cells
        in the rectangle $[x_1, x_2] \times [y_1, y_2]$.
\end{enumerate}

\section{Solution: 2D Binary Indexed Tree}

A 2D BIT (Fenwick tree) extends the classical 1D structure to support
point updates and prefix-sum queries in two dimensions.

\paragraph{Update.} To add $v$ to position $(x, y)$ (1-indexed):
\begin{lstlisting}[numbers=none,frame=none]
for (int i = x; i <= S; i += i & (-i))
    for (int j = y; j <= S; j += j & (-j))
        tree[i][j] += v;
\end{lstlisting}

\paragraph{Prefix sum.}
$\displaystyle\mathrm{sum}(x, y) = \sum_{i=1}^{x}\sum_{j=1}^{y} a[i][j]$
is computed by:
\begin{lstlisting}[numbers=none,frame=none]
int s = 0;
for (int i = x; i > 0; i -= i & (-i))
    for (int j = y; j > 0; j -= j & (-j))
        s += tree[i][j];
\end{lstlisting}

\paragraph{Rectangle sum.} By inclusion--exclusion:
\[
  \mathrm{RectSum}(x_1,y_1,x_2,y_2)
  = \mathrm{sum}(x_2,y_2)
  - \mathrm{sum}(x_1{-}1,y_2)
  - \mathrm{sum}(x_2,y_1{-}1)
  + \mathrm{sum}(x_1{-}1,y_1{-}1).
\]

\section{C++ Implementation}

\begin{lstlisting}
#include <bits/stdc++.h>
using namespace std;

int S;
int tree[1025][1025];

void update(int x, int y, int v) {
    for (int i = x; i <= S; i += i & (-i))
        for (int j = y; j <= S; j += j & (-j))
            tree[i][j] += v;
}

int query(int x, int y) {
    int s = 0;
    for (int i = x; i > 0; i -= i & (-i))
        for (int j = y; j > 0; j -= j & (-j))
            s += tree[i][j];
    return s;
}

int rectQuery(int x1, int y1, int x2, int y2) {
    return query(x2, y2)
         - query(x1 - 1, y2)
         - query(x2, y1 - 1)
         + query(x1 - 1, y1 - 1);
}

int main() {
    ios::sync_with_stdio(false);
    cin.tie(nullptr);

    int op;
    cin >> op >> S; // op == 0: initialise

    memset(tree, 0, sizeof(tree));

    while (cin >> op) {
        if (op == 3) break; // terminate

        if (op == 1) {
            int x, y, v;
            cin >> x >> y >> v;
            x++; y++; // convert 0-indexed input to 1-indexed BIT
            update(x, y, v);
        } else if (op == 2) {
            int l, b, r, t;
            cin >> l >> b >> r >> t;
            l++; b++; r++; t++; // convert to 1-indexed
            cout << rectQuery(l, b, r, t) << "\n";
        }
    }

    return 0;
}
\end{lstlisting}

\section{Complexity Analysis}
\begin{itemize}
  \item \textbf{Update:} $O(\log^2 S)$.
  \item \textbf{Query:} $O(\log^2 S)$ (four prefix-sum queries).
  \item \textbf{Space:} $O(S^2) = O(1024^2) \approx 10^6$.
  \item \textbf{Total time:} $O(Q \log^2 S)$ for $Q$ operations.
\end{itemize}

\end{document}
